\section{Algorithms}

\subsubsection{Sort and search}
\textbf{Search and sort from C}
\begin{itemize}
	\item \texttt{void *bsearch(const void *key, const void *base, \\
	size\_t n, size\_t size, \\
	int (*cmp) (const void *keyal, const void *datum))}\\
	Searches \texttt{base[0]} $\cdots$ \texttt{base[n-1]} for \texttt{key}.\\
	Comparison function must return negative if its first argument (search key) is less than its second argument (a table entry) and so on. 

	\item \texttt{qsort(void *base, size\_t n, size\_t size, \\
	int (*cmp)(const *void, const *void))}
\end{itemize}
\textbf{Search and sort from C++}\\

\begin{mdframed}[backgroundcolor=blue!20]
\textbf{Compare functions:}
\vspace{-8pt}
\begin{verbatim}
	bool compare(x, y){
	    return(x strictly precedes y)
	}
\end{verbatim}	
\end{mdframed}

In all the below, a could be an iterator or a pointer and all of them can use custom  compare functions.\\

\begin{itemize}
\item y = \texttt{lower\_bound(a, a+n, x)}\\
Returns an iterator to the first element whose value is $\geq$ x.\\
\item y = \texttt{upper\_bound(a, a+n, x)}\\
Returns an iterator to the first element whose value is $>$ x.\\
\item y = \texttt{equal\_range(a, a+n, x)}\\
Returns a pair of iterators pointing to the upper bound and lower bound of x.\\
\item y = \texttt{equal\_range(a, a+n, x, compare\_function)}\\
Returns a pair of iterators pointing to the upper bound and lower bound of x.\\
\item \texttt{sort(it1, it2, compare\_function);}
\item \texttt{stable\_sort(it1, it2, compare\_function);}
\item \texttt{is\_sorted(it1, it2, compare\_function);}
\item \texttt{is\_sorted\_until(it1, it2, compare\_function);}\\
Returns the iterator to the element until which the array is sorted.\\
\item \texttt{find\_if(start\_it, end\_it, func)}\\
Returns iterator to the first element in the range \texttt{[start\_it, end\_it)} for which \texttt{func} returns true.
\item \texttt{find\_if\_not(start\_it, end\_it, func)}\\
Returns iterator to the first element in the range \texttt{[start\_it, end\_it)} for which \texttt{func} returns false.
\end{itemize}


\subsubsection{Min Max}
\begin{itemize}
\item \texttt{max(a, b, comp);} \texttt{min(a, b, comp);}
\item \texttt{max(\{a1, a2, a3\}, comp);} \texttt{min(\{a1, a2, a3\}, comp);}
\item \texttt{max\_element(a, a+n, comp);}\\
Returns the iterator for the largest element in a. If there are multiple largest element, returns the iterator/pointer for the first one.
\item \texttt{min\_element(a, a+n, comp);}\\
\end{itemize}

\subsubsection{Merge}
\begin{itemize}
\item \texttt{merge(InputIterator F1, InputIterator L1, InputIterator F2, InputIterator L2, OutputIterator R);} \\
Returns the iterator to the end of output.\\
NOTE: This operator does not allocate memory. The \texttt{OutputIterator R} should point to a vector with sufficient memory to store the concatenated vector.\\
\item \texttt{set\_intersection} Arguments and output are the same as that for \texttt{merge}.
\item \texttt{set\_union} Arguments and output are the same as that for \texttt{merge}.
\item \texttt{set\_difference} Arguments and output are the same as that for \texttt{merge}.
\item \texttt{set\_symmetric\_difference} Arguments and output are the same as that for \texttt{merge}.
\end{itemize}

\subsubsection{Hash Function}
Can be used for custom types in \texttt{unordered\_map} and \texttt{unordered\_set}.
\begin{mdframed}[backgroundcolor=gray!10,linecolor=Firebrick4]
Usage:\\
\texttt{hash<T> hasher;}\\
\texttt{size\_t hasvalue = hasher(x);//Where x is an object of type T}\\
\end{mdframed}
Eg:
\begin{verbatim}
    hash<int> h; size_t x = h(10);
    hash<string> h; size_t x = h("hello");
\end{verbatim}

\textbf{Custom hash function.}

% \textbf{Custom hash function.}
Custom hash function for a user defined type can be defined as follows:
\begin{verbatim}
   class CustomHash {
    public:
    size_t operator () (const pair<string, string>& p) const {
        hash<string> hasher;
        auto h1 = hasher(p.first);
        auto h2 = hasher(p.second);
        return h1 ^ h2; // Below for better heuristics.
    }
}; 
\end{verbatim}

\textbf{Heuristics to combine hash values.}\\
\texttt{h2 = h2 \textasciicircum (h1 << 1);}
\texttt{h2 \textasciicircum= h2 + 0x9e3779b9 + (h1<<6) + (h1>>2);}\\
0x9e3779b9 is a prime number.\\
See: https://stackoverflow.com/a/2595226/5607735\\


\subsubsection{Others}

\begin{itemize}
\item \texttt{swap(T\& a, T\&b);}\\
\item \texttt{reverse(Iterator first, Iterator last);}
\item \texttt{next\_permutation(a.begin(), a.end());}\\
\item \texttt{prev\_permutation(a.begin(), a.end());}\\
\item \texttt{random\_shuffle(a.begin(), a.end());}\\
Gives the next permutation for a given vector of integers.
\end{itemize}

\vfill \null
\columnbreak 

